\documentclass{article}
\usepackage{amsmath,hyperref}
\title{A finite common-profile comparison for the warehouse receiver}
\date{}
\begin{document}
\maketitle

\paragraph{Result and scope.} In the backed primary model, compare a tick-$-1$ issue with first-hop onset ahead of due delivery (source) against a tick-0 issue without the hop edit (target). These are not identical histories: the source packet was in flight before onset and the target packet had not yet been issued. Both bring the named order to support 1 at tick 1. Both live cuts, both terminal dwell settings and deadlines 4--6 are admitted. The target includes both orders of channel onset processing, even though with no hop edit these have identical interface effects. This verifies one receiver-level comparison; it neither models arbitrary warehouse dynamics nor selects a held-out transfer.

\paragraph{Native domains and maps.} At every simulated pre-start phase reconstruct from actual placement, acceptance, begin and receipt events the witness of the channel to support 1, the witness of adjacent supports, and the witness of the local clock. Verify active-authorization guard snapshots against the controller's own local log. Domain projections discard unneeded global events but keep native provenance; in particular the channel domain includes the actual issue tick, and the variant is retained even when no receiver can distinguish it. Witnesses carry the shared external cut key $(\text{settlement prefix},\text{terminal dwell edit},\text{deadline},\text{tick})$. This is a synchronization index for the verifier, not a controller observation. Each map sends a source witness to the target's ordinary-A witness with that key. The ordinary-B witnesses are not in the raw image. These are not maps on complete packet traces, or an assertion of equivalence of the two plants in all external contexts.

The primitive factors are the received named authorization history (including receipt tick), adjacent load and reservation reads including terminal release as visible here, and clock and deadline; cut equality across channel/support and support/clock; and the receiver's own start decision, undefined for a mismatched raw tuple. Compatible tuples are a receiver observation assembly and may contain representatives with different provenance, not evidence of an executable global trace. Actual reachable prefixes are checked independently against simulated traces. The physical primitive excludes the root-acceptance event whose derivation is sought; adjacent stocks after earlier transfers can still be inspected by the receiver.

\paragraph{Reflection and coverage.} The pinned theorem \href{https://github.com/eternal-november/eos-research-programmes/blob/e1348964f73704212d2d8b48d5982a5ed600a2ff/projects/contextual-composition/feedback-and-concurrency/observable-covers/Residual_Profiles_and_Admission_Reflection.tex}{Existence and construction of common profiles} requires both test reflection T and target residual-class coverage C. For every component the maps preserve its key and its output payload. Therefore every unary read, cut synchronization and joint start test is reflected on the full raw source product, including incompatible raw tuples rather than only matching executions. There are 884,736 raw source triples and 884,640 are incompatible; the factors also check 18,432 raw binary pairs and 288 unary rows. The exact target residual classes are key and output payload: distinct payloads are separated by the unary read, distinct keys by a matching-key witness in a neighboring component; equal key and payload preserve every factor. Each target domain has 192 native witnesses, with 96 residual classes hit by the 96 source witnesses. This verifies C separately from T, without requiring native surjectivity. A diagnostic comparison with no hop edit on the early path fails immediately: at tick 0 the target has the named permit and has a new authorization class. It is not an independently chosen rejection case.

\paragraph{Through-prefix receiver guarantee.} Initially no root acceptance is pending. At a matching phase with the same residual observations, the common factor evaluation transports the start decision. A root start schedules the unique root acceptance exactly one tick later: the two endpoints remain reserved, the root keeps the load until acceptance, and no competing start can modify either endpoint. No other event in this plant can move the root's load. This residual grant relation is preserved from the empty initial grant state by the actual acceptance/start phase order. The projection separately replays both finite plants to check that all subsequent compared phases continue to have matching inputs (the theorem does not infer reachability or this input relation). Consequently the unchanged receiver makes the same choice at each matched phase, and every such start yields the same one-tick actual root release and service-by-deadline result. This is conditional on the primitive contracts and checked interface histories; it is not a circular assumption of equal end-to-end service or transport of a whole-plant safety theorem.

\paragraph{Matched direct comparison.} Both proofs need the same primary replay/prefix laws, native receipt-time derivation, adjacent resource/reservation evaluation, local clock data, unchanged guard, and one-tick pending-grant invariant. The direct proof combines those matched inputs in the cached guard and then invokes the same pending-grant invariant. The profile proof additionally fixes the native maps, checks raw factors, and proves residual coverage. The direct argument does not need coverage for its pointwise transfer claim, so there is no demonstrated reduction of proof work. Twelve source scenario replays and twenty-four target replays are shared; twelve extra local diagnostic replays do not count as held-out validation. No timing or labor savings are inferred.

\paragraph{Proposed transfer admissibility and freeze.} Freeze the input/output tests, code, map discipline (compare cases using the actual alignment of the external phase, without rewriting observational values), root grant invariant and cached direct baseline before a non-author chooses a source extension with a genuinely changed cause of permit delivery or adjacent physical availability. The plant must still provide independently executed local histories, identified physical endpoints and resource law, atomic observations with the same declared meanings, and next-tick completion of the root grant after valid dispatch. Class coverage, raw reflection, projection along the new prefixes and a truly unchanged receiver are review tests, not assumptions that reject an inconvenient result up front. The non-author specifies the extension, coupling and scenario support without conditioning on which tests pass. A positive transfer must present actual substituted maps and discharge all tests; a deliberately rejected transfer must exhibit a failed primitive equality or uncovered residual class. Keep failures visible; a subsequent interface revision requires re-freezing and a new independent case. Merely rerunning or relabeling the current cases does not qualify.
\end{document}
